\chapter{Mouvement circulaire et modèles sinusoïdaux}
\label{workbook:chapter:17}
\section{Vitesse angulaire}
\label{workbook:section:17.01}
\begin{exerciseblock}{Vitesse angulaire}
\item Une roue fait $18$ tours par minute. Calculer sa vitesse angulaire en rad/s.
\item Avec $\omega=5$ rad/s, calculer la période et la fréquence.
\item Comparer vitesse linéaire au bord de deux roues de rayons différents tournant avec la même $\omega$.
\end{exerciseblock}
\section{Coordonnées d'un mouvement circulaire}
\label{workbook:section:17.02}
\begin{exerciseblock}{Coordonnées d'un mouvement circulaire}
\item Écrire les coordonnées d'un point de rayon $6$, centre $(2,-1)$, angle initial $\pi/3$ et vitesse angulaire $2$.
\item Calculer sa position à $t=\pi/4$.
\item Montrer que la trajectoire satisfait $(x-2)^2+(y+1)^2=36$.
\end{exerciseblock}
\section{Construire un modèle à partir de données}
\label{workbook:section:17.03}
\begin{exerciseblock}{Construire un modèle à partir de données}
\item Une grande roue a hauteur minimale $2$ m, maximale $42$ m et période $90$ s. Trouver amplitude, médiane et fréquence angulaire.
\item Si la nacelle part au point le plus bas, construire un modèle cosinus.
\item Si elle part sur la médiane en montant, construire un modèle sinus.
\end{exerciseblock}
\section{Interpréter un modèle}
\label{workbook:section:17.04}
\begin{exerciseblock}{Interpréter un modèle}
\item Pour $h(t)=12+7\sin(\pi(t-3)/4)$, interpréter chaque paramètre.
\item Trouver les premiers instants où $h(t)=12$, le maximum et le minimum.
\item Déterminer un domaine temporel pertinent si le modèle décrit seulement trois cycles observés.
\end{exerciseblock}
\section{Interprétation ondulatoire}
\label{workbook:section:17.05}
\begin{exerciseblock}{Interprétation ondulatoire}
\item Pour $y(x,t)=3\sin(2\pi x/5-4\pi t)$, identifier amplitude, longueur d'onde, fréquence et vitesse de propagation.
\item Comparer l'effet d'une amplitude doublée et d'une fréquence doublée sur une onde sonore idéale.
\item Expliquer la différence entre période temporelle et longueur d'onde spatiale.
\end{exerciseblock}
\section{Du mouvement circulaire au modèle temporel}
\label{workbook:section:17.06}
\begin{exerciseblock}{Du mouvement circulaire au modèle temporel}
\item Un point tourne avec rayon $r$, hauteur du centre $D$, vitesse $\omega$ et phase $\phi$. Écrire sa hauteur générale.
\item Déduire maximum, minimum et période sans résoudre d'équation.
\item Expliquer comment choisir sinus ou cosinus selon la position initiale.
\end{exerciseblock}
\section{Construire un modèle périodique à partir d'observations}
\label{workbook:section:17.07}
\begin{exerciseblock}{Construire un modèle périodique à partir d'observations}
\item Des maxima sont observés aux temps $2$, $10$, $18$ avec valeurs $15$; les minima valent $3$. Construire un modèle.
\item Un phénomène traverse sa médiane $7$ en montant à $t=1$ et a période $12$, amplitude $4$. Écrire un modèle.
\item Décrire comment détecter une phase incorrecte en comparant le sens de variation initial.
\end{exerciseblock}
\section{Construire et valider un modèle périodique}
\label{workbook:section:17.08}
\begin{exerciseblock}{Construire et valider un modèle périodique}
\item À partir de données bruitées autour d'une oscillation, proposer une méthode pour estimer médiane, amplitude et période.
\item Valider le modèle $T(t)=18+5\cos(2\pi(t-4)/24)$ pour une température journalière : interpréter et vérifier les extrêmes.
\item Donner trois raisons pour lesquelles un phénomène périodique réel peut ne pas être exactement sinusoïdal.
\end{exerciseblock}
\section*{Synthèse du chapitre}
\begin{synthesisbox}
\item Une marée a maximum $4{,}8$ m à $3$ h et minimum $0{,}8$ m à $9{,}2$ h. Supposer un comportement sinusoïdal : construire un modèle et prédire la hauteur à $6$ h.
\item Une roue de rayon $0{,}35$ m tourne à $120$ tr/min. Calculer $\omega$, la période et la vitesse linéaire d'un point du bord.
\item Un modèle périodique prédit correctement les extrema mais les passages par la médiane sont systématiquement trop tôt. Quel paramètre faut-il revoir en priorité? Justifier.
\end{synthesisbox}

